\documentclass[10pt,a4paper]{amsart}
\usepackage{amsmath,amssymb,amsthm,graphicx}
\usepackage{verbatim}
\newenvironment{reference}{\comment}{\endcomment}
%\newenvironment{reference}{}{}
\newenvironment{slogan}{\comment}{\endcomment}
\newenvironment{history}{\comment}{\endcomment}
% For commutative diagrams we use Xy-pic
\usepackage[all]{xy}

% We use 2cell for 2-commutative diagrams.
\xyoption{2cell}
% Theorem environments.
%
\theoremstyle{plain}
\newtheorem{theorem}[subsection]{Theorem}
\newtheorem{proposition}[subsection]{Proposition}
\newtheorem{lemma}[subsection]{Lemma}

\theoremstyle{definition}
\newtheorem{definition}[subsection]{Definition}
\newtheorem{example}[subsection]{Example}
\newtheorem{exercise}[subsection]{Exercise}
\newtheorem{situation}[subsection]{Situation}

\theoremstyle{remark}
\newtheorem{remark}[subsection]{Remark}
\newtheorem{remarks}[subsection]{Remarks}

\numberwithin{equation}{subsection}

% Macros
%
\def\lim{\mathop{\mathrm{lim}}\nolimits}
\def\colim{\mathop{\mathrm{colim}}\nolimits}
\def\Spec{\mathop{\mathrm{Spec}}}
\def\Hom{\mathop{\mathrm{Hom}}\nolimits}
\def\Ext{\mathop{\mathrm{Ext}}\nolimits}
\def\SheafHom{\mathop{\mathcal{H}\!\mathit{om}}\nolimits}
\def\SheafExt{\mathop{\mathcal{E}\!\mathit{xt}}\nolimits}
\def\Sch{\mathit{Sch}}
\def\Mor{\mathop{\mathrm{Mor}}\nolimits}
\def\Ob{\mathop{\mathrm{Ob}}\nolimits}
\def\Sh{\mathop{\mathit{Sh}}\nolimits}
\def\NL{\mathop{N\!L}\nolimits}
\def\CH{\mathop{\mathrm{CH}}\nolimits}
\def\proetale{{pro\text{-}\acute{e}tale}}
\def\etale{{\acute{e}tale}}
\def\QCoh{\mathit{QCoh}}
\def\Ker{\mathop{\mathrm{Ker}}}
\def\Im{\mathop{\mathrm{Im}}}
\def\Coker{\mathop{\mathrm{Coker}}}
\def\Coim{\mathop{\mathrm{Coim}}}

% Boxtimes
%
\DeclareMathSymbol{\boxtimes}{\mathbin}{AMSa}{"02}

%
% Macros for moduli stacks/spaces
%
\def\QCohstack{\mathcal{QC}\!\mathit{oh}}
\def\Cohstack{\mathcal{C}\!\mathit{oh}}
\def\Spacesstack{\mathcal{S}\!\mathit{paces}}
\def\Quotfunctor{\mathrm{Quot}}
\def\Hilbfunctor{\mathrm{Hilb}}
\def\Curvesstack{\mathcal{C}\!\mathit{urves}}
\def\Polarizedstack{\mathcal{P}\!\mathit{olarized}}
\def\Complexesstack{\mathcal{C}\!\mathit{omplexes}}
% \Pic is the operator that assigns to X its picard group, usage \Pic(X)
% \Picardstack_{X/B} denotes the Picard stack of X over B
% \Picardfunctor_{X/B} denotes the Picard functor of X over B
\def\Pic{\mathop{\mathrm{Pic}}\nolimits}
\def\Picardstack{\mathcal{P}\!\mathit{ic}}
\def\Picardfunctor{\mathrm{Pic}}
\def\Deformationcategory{\mathcal{D}\!\mathit{ef}}
\usepackage{cleveref}
\crefname{originallemma}{original lemma}{original lemmas}
\Crefname{originallemma}{Original lemma}{Original lemmas}
\makeatletter
\newlabel{sites-cohomology-lemma-relative-Leray}{{0734}{source}{}{sourcepointer-1}{}}
\newlabel{sites-cohomology-lemma-relative-Leray@cref}{{[originallemma][0][]0734}{[1][0][]source}{}{}{}}
\makeatother
\crefname{originallemma}{original lemma}{original lemmas}
\Crefname{originallemma}{Original lemma}{Original lemmas}
\makeatletter
\newlabel{spaces-morphisms-lemma-pushforward}{{03M9}{source}{}{sourcepointer-2}{}}
\newlabel{spaces-morphisms-lemma-pushforward@cref}{{[originallemma][0][]03M9}{[1][0][]source}{}{}{}}
\makeatother
\crefname{originallemma}{original lemma}{original lemmas}
\Crefname{originallemma}{Original lemma}{Original lemmas}
\makeatletter
\newlabel{spaces-properties-lemma-characterize-quasi-coherent}{{03M0}{source}{}{sourcepointer-3}{}}
\newlabel{spaces-properties-lemma-characterize-quasi-coherent@cref}{{[originallemma][0][]03M0}{[1][0][]source}{}{}{}}
\makeatother
\crefname{originallemma}{original lemma}{original lemmas}
\Crefname{originallemma}{Original lemma}{Original lemmas}
\makeatletter
\newlabel{spaces-properties-lemma-properties-quasi-coherent}{{03M1}{source}{}{sourcepointer-4}{}}
\newlabel{spaces-properties-lemma-properties-quasi-coherent@cref}{{[originallemma][0][]03M1}{[1][0][]source}{}{}{}}
\makeatother
\crefname{originallemma}{original lemma}{original lemmas}
\Crefname{originallemma}{Original lemma}{Original lemmas}
\makeatletter
\newlabel{spaces-properties-lemma-pushforward-etale-base-change-modules}{{03LX}{source}{}{sourcepointer-5}{}}
\newlabel{spaces-properties-lemma-pushforward-etale-base-change-modules@cref}{{[originallemma][0][]03LX}{[1][0][]source}{}{}{}}
\makeatother
\crefname{originallemma}{original lemma}{original lemmas}
\Crefname{originallemma}{Original lemma}{Original lemmas}
\makeatletter
\newlabel{spaces-properties-lemma-quasi-compact-affine-cover}{{03H6}{source}{}{sourcepointer-6}{}}
\newlabel{spaces-properties-lemma-quasi-compact-affine-cover@cref}{{[originallemma][0][]03H6}{[1][0][]source}{}{}{}}
\makeatother
\title{Higher direct images of a quasi-coherent module along a quasi-compact quasi-separated morphism of algebraic spaces are quasi-coherent}
\author{The Stacks Project Authors}
\date{stacks-project:spaces-cohomology.tex:local-snapshot-20261012; GNU FDL 1.2 or later}
\begin{document}
\maketitle
\noindent\textit{Editorial note.} All delivered mathematical text is the author's own, copied from the pinned Stacks Project chapter file spaces-cohomology.tex; no statement or proof step was rewritten, reordered or removed.\par\medskip
\noindent\textit{Editorial note.} Preamble provenance: the pinned chapter file opens with \textbackslash{}input\{preamble\} and carries no preamble text of its own, so the project's single Stacks-wide preamble staged at sources/stacks-spaces-cohomology/preamble.tex is excerpted verbatim as three spans, lines 8-12 (verbatim package and the project's own reference/slogan/history display environments), lines 14-18 (Xy-pic setup) and lines 154-217 (theorem declarations and mathematical macros). The wrapper class and the amsmath/amssymb/amsthm/graphicx base packages are supplied by the assembler. No macro was invented or redefined.\par\medskip
\noindent\textit{Editorial note.} Same-chapter pointers are located at their \textbackslash{}label lines inside the pinned chapter file; cross-chapter pointers are located in the project's own tag list staged in the snapshot (sources/stacks-chow/tags.txt), at the verbatim lines carrying their canonical tags.\par\medskip
\noindent\textit{Allowed standard prerequisites (editorial).}\begin{itemize}
\item Base change of higher direct images along an etale localization, and etale locality of quasi-coherence. Properties of Spaces, labels spaces-properties-lemma-pushforward-etale-base-change-modules (invoked at proof line 351) and spaces-properties-lemma-characterize-quasi-coherent (invoked at proof line 357). Together they reduce the claim to the case where the target \$Y\$ is an affine scheme.
\item An affine etale cover of a quasi-compact algebraic space. Properties of Spaces, label spaces-properties-lemma-quasi-compact-affine-cover (invoked at proof line 365). It provides the surjective etale morphism \$g : U \textbackslash{}to X\$ with \$U\$ affine, used to run the induction through the Leray spectral sequence of \$f \textbackslash{}circ g\$.
\item Quasi-coherence of the zeroth direct image. Morphisms of Spaces, label spaces-morphisms-lemma-pushforward (invoked at proof line 385). It is the base case \$n = 0\$ of the induction \$IH\_n\$.
\item The Leray spectral sequence of a composition. Cohomology on Sites, label sites-cohomology-lemma-relative-Leray (invoked at proof line 396). It relates the higher direct images of \$f \textbackslash{}circ g\$ to those of \$f\$ and \$g\$, which is how the induction step is organized.
\item Quasi-coherent modules form a weak Serre subcategory of all modules. Properties of Spaces, label spaces-properties-lemma-properties-quasi-coherent (invoked at proof line 408). It is what upgrades the information from the spectral sequence, all of whose other terms in the relevant degrees are known quasi-coherent, into quasi-coherence of \$R\textasciicircum{}\{n+1\}f\_*g\_*\textbackslash{}mathcal\{H\}\$.
\end{itemize}\medskip
\ifdefined\hypertarget\hypertarget{sourcepointer-1}{}\fi
\noindent\textit{Source reference.} Original lemma 0734: sources/stacks-chow/tags.txt, source lines 9048--9048. Navigation pointer for \textbackslash{}ref\{sites-cohomology-lemma-relative-Leray\} used inside the counted span (line 396); the argument that consumes it is retained in the span.\par\medskip
\ifdefined\hypertarget\hypertarget{sourcepointer-2}{}\fi
\noindent\textit{Source reference.} Original lemma 03M9: sources/stacks-chow/tags.txt, source lines 4741--4741. Navigation pointer for \textbackslash{}ref\{spaces-morphisms-lemma-pushforward\} used inside the counted span (line 385); the argument that consumes it is retained in the span.\par\medskip
\ifdefined\hypertarget\hypertarget{sourcepointer-3}{}\fi
\noindent\textit{Source reference.} Original lemma 03M0: sources/stacks-chow/tags.txt, source lines 4732--4732. Navigation pointer for \textbackslash{}ref\{spaces-properties-lemma-characterize-quasi-coherent\} used inside the counted span (line 357); the argument that consumes it is retained in the span.\par\medskip
\ifdefined\hypertarget\hypertarget{sourcepointer-4}{}\fi
\noindent\textit{Source reference.} Original lemma 03M1: sources/stacks-chow/tags.txt, source lines 4733--4733. Navigation pointer for \textbackslash{}ref\{spaces-properties-lemma-properties-quasi-coherent\} used inside the counted span (line 408); the argument that consumes it is retained in the span.\par\medskip
\ifdefined\hypertarget\hypertarget{sourcepointer-5}{}\fi
\noindent\textit{Source reference.} Original lemma 03LX: sources/stacks-chow/tags.txt, source lines 4729--4729. Navigation pointer for \textbackslash{}ref\{spaces-properties-lemma-pushforward-etale-base-change-modules\} used inside the counted span (line 351); the argument that consumes it is retained in the span.\par\medskip
\ifdefined\hypertarget\hypertarget{sourcepointer-6}{}\fi
\noindent\textit{Source reference.} Original lemma 03H6: sources/stacks-chow/tags.txt, source lines 4555--4555. Navigation pointer for \textbackslash{}ref\{spaces-properties-lemma-quasi-compact-affine-cover\} used inside the counted span (line 365); the argument that consumes it is retained in the span.\par\medskip
\begin{lemma}
\label{lemma-higher-direct-image}
Let $S$ be a scheme. Let $f : X \to Y$ be a morphism of algebraic spaces
over $S$. If $f$ is quasi-compact and quasi-separated, then $R^if_*$
transforms quasi-coherent $\mathcal{O}_X$-modules into
quasi-coherent $\mathcal{O}_Y$-modules.
\end{lemma}

\begin{proof}
Let $V \to Y$ be an \'etale morphism where $V$ is an affine scheme. Set
$U = V \times_Y X$ and denote $f' : U \to V$ the induced morphism.
Let $\mathcal{F}$ be a quasi-coherent $\mathcal{O}_X$-module. By
Properties of Spaces, Lemma
\ref{spaces-properties-lemma-pushforward-etale-base-change-modules}
we have
$R^if'_*(\mathcal{F}|_U) = (R^if_*\mathcal{F})|_V$.
Since the property of being a quasi-coherent module is local in the
\'etale topology on $Y$ (see
Properties of Spaces, Lemma
\ref{spaces-properties-lemma-characterize-quasi-coherent})
we may replace $Y$ by $V$, i.e., we may assume $Y$ is an affine scheme.

\medskip\noindent
Assume $Y$ is affine. Since $f$ is quasi-compact we see that $X$
is quasi-compact. Thus we may choose an affine scheme $U$ and a surjective
\'etale morphism $g : U \to X$, see
Properties of Spaces,
Lemma \ref{spaces-properties-lemma-quasi-compact-affine-cover}.
Picture
$$
\xymatrix{
U \ar[r]_g \ar[rd]_{f \circ g} & X \ar[d]^f \\
& Y
}
$$
The morphism $g : U \to X$ is representable, separated
and quasi-compact because $X$ is quasi-separated. Hence the lemma
holds for $g$ (by the discussion above the lemma).
It also holds for $f \circ g : U \to Y$ (as this is a morphism
of affine schemes).

\medskip\noindent
In the situation described in the previous paragraph we will show by
induction on $n$ that $IH_n$: for any quasi-coherent sheaf $\mathcal{F}$
on $X$ the sheaves $R^if\mathcal{F}$
are quasi-coherent for $i \leq n$.
The case $n = 0$ follows from
Morphisms of Spaces, Lemma \ref{spaces-morphisms-lemma-pushforward}.
Assume $IH_n$. In the rest of the proof we show that $IH_{n + 1}$ holds.

\medskip\noindent
Let $\mathcal{H}$ be a quasi-coherent $\mathcal{O}_U$-module.
Consider the Leray spectral sequence
$$
E_2^{p, q} = R^pf_* R^qg_* \mathcal{H}
\Rightarrow
R^{p + q}(f \circ g)_*\mathcal{H}
$$
Cohomology on Sites, Lemma \ref{sites-cohomology-lemma-relative-Leray}.
As $R^qg_*\mathcal{H}$ is quasi-coherent by $IH_n$ all the sheaves
$R^pf_*R^qg_*\mathcal{H}$ are quasi-coherent for $p \leq n$.
The sheaves $R^{p + q}(f \circ g)_*\mathcal{H}$ are all
quasi-coherent (in fact zero for $p + q > 0$ but we do not need this).
Looking in degrees $\leq n + 1$ the only module which we do not
yet know is quasi-coherent is $E_2^{n + 1, 0} = R^{n + 1}f_*g_*\mathcal{H}$.
Moreover, the differentials
$d_r^{n + 1, 0} : E_r^{n + 1, 0} \to E_r^{n + 1 + r, 1 - r}$
are zero as the target is zero. Using that $\QCoh(\mathcal{O}_X)$
is a weak Serre subcategory of $\textit{Mod}(\mathcal{O}_X)$
(Properties of Spaces, Lemma
\ref{spaces-properties-lemma-properties-quasi-coherent}) it
follows that $R^{n + 1}f_*g_*\mathcal{H}$
is quasi-coherent (details omitted).

\medskip\noindent
Let $\mathcal{F}$ be a quasi-coherent $\mathcal{O}_X$-module.
Set $\mathcal{H} = g^*\mathcal{F}$. The adjunction mapping
$\mathcal{F} \to g_*g^*\mathcal{F} = g_*\mathcal{H}$ is injective
as $U \to X$ is surjective \'etale. Consider the exact sequence
$$
0 \to \mathcal{F} \to g_*\mathcal{H} \to \mathcal{G} \to 0
$$
where $\mathcal{G}$ is the cokernel of the first map and in particular
quasi-coherent. Applying the long exact cohomology sequence we obtain
$$
R^nf_*g_*\mathcal{H} \to
R^nf_*\mathcal{G} \to
R^{n + 1}f_*\mathcal{F} \to
R^{n + 1}f_*g_*\mathcal{H} \to
R^{n + 1}f_*\mathcal{G}
$$
The cokernel of the first arrow is quasi-coherent and
we have seen above that $R^{n + 1}f_*g_*\mathcal{H}$ is quasi-coherent.
Thus $R^{n + 1}f_*\mathcal{F}$ has a $2$-step filtration
where the first step is quasi-coherent and the second a submodule of
a quasi-coherent sheaf. Since $\mathcal{F}$ is an arbitrary quasi-coherent
$\mathcal{O}_X$-module, this result also holds for $\mathcal{G}$.
Thus we can choose an exact sequence
$0 \to \mathcal{A} \to R^{n + 1}f_*\mathcal{G} \to \mathcal{B}$
with $\mathcal{A}$, $\mathcal{B}$ quasi-coherent $\mathcal{O}_Y$-modules.
Then the kernel $\mathcal{K}$ of
$R^{n + 1}f_*g_*\mathcal{H} \to R^{n + 1}f_*\mathcal{G}
\to \mathcal{B}$ is quasi-coherent, whereupon we obtain a map
$\mathcal{K} \to \mathcal{A}$ whose kernel $\mathcal{K}'$ is
quasi-coherent too. Hence $R^{n + 1}f_*\mathcal{F}$ sits in an exact
sequence
$$
R^nf_*g_*\mathcal{H} \to
R^nf_*\mathcal{G} \to
R^{n + 1}f_*\mathcal{F} \to \mathcal{K}' \to 0
$$
with all modules quasi-coherent except for possibly $R^{n + 1}f_*\mathcal{F}$.
We conclude that $R^{n + 1}f_*\mathcal{F}$ is quasi-coherent, i.e.,
$IH_{n + 1}$ holds as desired.
\end{proof}

\end{document}
